\documentclass[11pt]{amsart}

\usepackage[T1]{fontenc}
\usepackage{lmodern}
\usepackage{microtype}
\usepackage{amsmath,amssymb,amsthm}
\usepackage[hidelinks]{hyperref}

\hypersetup{
  pdftitle={
    Counterexamples to Gyoda's Injectivity Conjecture
    for Generalized Markov Trees
  }
}

\newtheorem{theorem}{Theorem}
\newtheorem{lemma}[theorem]{Lemma}

\newcommand{\GMT}{\mathrm{M}\mathbb{T}}
\newcommand{\id}{\mathrm{id}}

\title[Gyoda's injectivity conjecture]{
  Counterexamples to Gyoda's Injectivity Conjecture
  for Generalized Markov Trees
}
\date{}

\subjclass[2020]{11D25, 11B39}
\keywords{
  generalized Markov equation, generalized Markov tree, Farey labeling,
  uniqueness conjecture, Fibonacci numbers}

\begin{document}
\raggedbottom

\begin{abstract}
Gyoda conjectured that, in each fixed generalized Markov tree, equality of
numerical labels $n_s=n_t$ forces equality of the corresponding Farey labels.
We refute this conjecture by the collision
\[
 n_{1/4}=n_{2/3}=889
 \quad\text{in}\quad
 \GMT(22,0,5,\id),
\]
whose parameter triple is pairwise distinct. We also construct infinitely
many counterexample trees, beginning with
$n_{1/5}=n_{2/3}=89$ in $\GMT(6,0,0,\id)$. In all of these examples the
equal numerical labels have different coordinate-position tags; thus the
result concerns the map $t\mapsto n_t$, exactly as in the stated conjecture.
\end{abstract}

\maketitle

\section{The conjecture and the result}

For $k_1,k_2,k_3\in\mathbb Z_{\ge 0}$, the generalized Markov equation is
\begin{equation}\label{eq:GM}
 x^2+y^2+z^2+k_1yz+k_2zx+k_3xy
 =(3+k_1+k_2+k_3)xyz.
\end{equation}
Gyoda and Matsushita introduced its mutation structure
in~\cite{GyodaMatsushita}. Gyoda subsequently synchronized the resulting
binary trees with the Farey tree and formulated the following assertion as
Conjecture~7.6 in~\cite{Gyoda}, where he recorded that no counterexample was
known: for every fixed tree $\GMT(k_1,k_2,k_3,\sigma)$ and irreducible
fractions $s,t\in[0,1]$,
\[
 n_s=n_t \quad\Longrightarrow\quad s=t.
\]
Here $(n_t,i_t)$ denotes the generalized Markov number and its coordinate
position at the node whose middle Farey label is $t$. A different
generalization of the uniqueness conjecture for the same
equation~\eqref{eq:GM}, stated for size-ordered triples rather than for the
labeling $t\mapsto n_t$, had been proposed shortly before by Chen and Jia
\cite[Conjecture~8.2]{ChenJia}; it is delimited from Conjecture~7.6 in the
Status and scope section below.

For completeness, we recall the part of the construction used below. The
Farey tree has root $(0/1,1/1,1/0)$, and the children of a Farey triple
$(r,t,s)$ are
\[
 (r,r\mathbin\oplus t,t)
 \quad\text{and}\quad
 (t,t\mathbin\oplus s,s),
 \qquad
 \frac ab\mathbin\oplus\frac cd:=\frac{a+c}{b+d}.
\]
For $\sigma=\id$, the corresponding generalized Markov tree has root
\[
 ((1,1),(k_2+2,2),(1,3)).
\]
If $v=((a,h),(b,i),(c,j))$, its left and right children are
\begin{align*}
 L(v)
 &=
 \left(
   (a,h),
   \left(\frac{a^2+k_jab+b^2}{c},j\right),
   (b,i)
 \right),\\
 R(v)
 &=
 \left(
   (b,i),
   \left(\frac{b^2+k_hbc+c^2}{a},h\right),
   (c,j)
 \right).
\end{align*}
The middle pair at the node corresponding to $(r,t,s)$ is $(n_t,i_t)$.

\begin{theorem}\label{thm:main}
Gyoda's Conjecture~7.6 is false.
\begin{enumerate}
\item
In $\GMT(22,0,5,\id)$,
\[
 (n_{1/4},i_{1/4})=(889,3),
 \qquad
 (n_{2/3},i_{2/3})=(889,1).
\]
In particular, the conjecture fails for a pairwise distinct parameter triple.

\item
Let $F_0=0$, $F_1=1$, and
$F_{r+1}=F_r+F_{r-1}$ for $r\ge1$. For every $m\ge0$, set
\[
 d_m=30m+5,
 \qquad
 q_m=\frac{F_{60m+11}-29}{10}.
\]
Then $q_m\in\mathbb Z_{\ge0}$ and, in $\GMT(q_m,0,0,\id)$,
\[
 (n_{1/d_m},i_{1/d_m})=(F_{60m+11},2),
 \qquad
 (n_{2/3},i_{2/3})=(F_{60m+11},1).
\]
For $m=0$, this is the collision
$n_{1/5}=n_{2/3}=89$ in $\GMT(6,0,0,\id)$.
\end{enumerate}
\end{theorem}

\section{A pairwise distinct counterexample}

Take $(k_1,k_2,k_3)=(22,0,5)$. Starting at the root, three left mutations
give
\begin{align*}
 ((1,1),(2,2),(1,3))
 &\xrightarrow{L}((1,1),(15,3),(2,2))\\
 &\xrightarrow{L}((1,1),(113,2),(15,3))\\
 &\xrightarrow{L}((1,1),(889,3),(113,2)),
\end{align*}
where
\begin{gather*}
 15=1^2+5(1)(2)+2^2,
 \qquad
 113=\frac{1^2+15^2}{2},\\
 889=\frac{1^2+5(1)(113)+113^2}{15}.
\end{gather*}
The middle Farey labels along this path are $1/2$, $1/3$, and $1/4$.
On the other hand, the right child of the $1/2$ node is
\begin{gather*}
 R\bigl(((1,1),(15,3),(2,2))\bigr)
 =((15,3),(889,1),(2,2)),\\
 889=15^2+22(15)(2)+2^2,
\end{gather*}
whose middle Farey label is $2/3$. Hence
\[
 (n_{1/4},i_{1/4})=(889,3),
 \qquad
 (n_{2/3},i_{2/3})=(889,1).
\]
The position tags correspond to the ordered solutions
$(x,y,z)=(1,113,889)$ and $(x,y,z)=(889,2,15)$ of
\[
 x^2+y^2+z^2+22yz+5xy=30xyz,
\]
as substitution verifies:
\begin{align*}
 1+113^2+889^2+22(113)(889)+5(1)(113)
 &=3013710=30(1)(113)(889),\\
 889^2+2^2+15^2+22(2)(15)+5(889)(2)
 &=800100=30(889)(2)(15).
\end{align*}
This proves Theorem~\ref{thm:main}(1).

One remark on the pair, used nowhere in the proofs. The two colliding labels lie
at different depths of the Farey tree, $1/4$ at depth $3$ and $2/3$ at depth
$2$, but they have the same weight, that is, the same sum of numerator and
denominator, $1+4=2+3=5$; and $\{1/4,2/3\}$ is the first pair of distinct
irreducible fractions in $[0,1]$ of equal weight. A search for collisions that
enumerates Farey labels by weight rather than by depth therefore reaches this
pair first.

\section{An infinite family}

\begin{lemma}\label{lem:spine}
Fix $q\in\mathbb Z_{\ge0}$. In $\GMT(q,0,0,\id)$,
\[
 n_{1/d}=F_{2d+1}\quad(d\ge1),
 \qquad
 n_{2/3}=29+10q.
\]
Moreover, $i_{1/d}=2$ for odd $d$ and $i_{1/d}=3$ for even $d$, while
$i_{2/3}=1$.
\end{lemma}

\begin{proof}
Write $u_d=n_{1/d}$ and $i_d=i_{1/d}$. Since $k_2=0$, the root is
\[
 ((1,1),(2,2),(1,3)),
\]
whose middle Farey label is $1/1$, so $u_1=2$ and $i_1=2$. Its left child
is
\begin{align*}
 L\bigl(((1,1),(2,2),(1,3))\bigr)
 &=\left((1,1),\left(\frac{1^2+k_3(1)(2)+2^2}{1},3\right),(2,2)\right)\\
 &=((1,1),(5,3),(2,2)),
\end{align*}
with middle Farey label $1/2$; hence $u_2=5$ and $i_2=3$.

Along the left spine the first entry stays $(a,h)=(1,1)$, so the node labeled
$1/d$ is $((1,1),(u_d,i_d),(u_{d-1},i_{d-1}))$ and its left child is
\[
 \left(
  (1,1),
  \left(\frac{1+k_{i_{d-1}}u_d+u_d^2}{u_{d-1}},\,i_{d-1}\right),
  (u_d,i_d)
 \right).
\]
The new entry therefore inherits the tag of the third component of its
parent, and all tags occurring along the spine lie in $\{2,3\}$, where
$k_2=k_3=0$. Hence
\[
 u_{d+1}=\frac{u_d^2+1}{u_{d-1}},
 \qquad
 i_{d+1}=i_{d-1}
 \qquad(d\ge2),
\]
so the position tags alternate: with $i_1=2$ and $i_2=3$ we get $i_{1/d}=2$
for odd $d$ and $i_{1/d}=3$ for even $d$. The Fibonacci identity
\[
 F_{2d-1}F_{2d+3}=F_{2d+1}^2+1,
\]
which is the case $n=2d+1$, $r=2$ of Catalan's identity
$F_n^2-F_{n-r}F_{n+r}=(-1)^{n-r}F_r^2$, together with
$u_1=2=F_3$ and $u_2=5=F_5$, therefore proves by induction that
$u_d=F_{2d+1}$.

The node labeled $1/2$ is $((1,1),(5,3),(2,2))$. Its right child is
\[
 \left((5,3),\left(\frac{5^2+k_1(5)(2)+2^2}{1},1\right),(2,2)\right)
 =((5,3),(29+10q,1),(2,2)),
\]
which proves the remaining assertions.
\end{proof}

\begin{proof}[Proof of Theorem~\ref{thm:main}(2)]
A direct iteration of the Fibonacci recurrence modulo $10$ gives
\[
 (F_{60},F_{61})\equiv(0,1)\pmod{10}.
\]
The recurrence then implies
$F_{r+60}\equiv F_r\pmod{10}$ for every $r\ge0$. Hence
\[
 F_{60m+11}\equiv F_{11}=89\equiv9\pmod{10},
\]
so $q_m$ is an integer. Since $F_{60m+11}\ge F_{11}=89$, one also has
$q_m\ge6$.

Because $2d_m+1=60m+11$, Lemma~\ref{lem:spine} yields
\[
 n_{1/d_m}
 =F_{60m+11}
 =29+10q_m
 =n_{2/3}.
\]
Furthermore, $d_m$ is odd, so the corresponding position tags are $2$ and
$1$. Finally, $d_m\ge5$, and hence $1/d_m\ne2/3$. Since $q_m$ is strictly
increasing, these are infinitely many distinct counterexample trees.
\end{proof}

\section*{Status and scope}

The equal numerical labels in Theorem~\ref{thm:main} have different position
tags. Thus the theorem refutes Conjecture~7.6 exactly as stated, namely the
injectivity of $t\mapsto n_t$; it does not address injectivity of
$t\mapsto(n_t,i_t)$. Part~(2) gives infinitely many counterexample trees,
not infinitely many collisions in one fixed tree. No claim is made about
global minimality, the classical case $(0,0,0)$, or the symmetric locus
$k_1=k_2=k_3$.

Theorem~\ref{thm:main} is also not a restatement of an already known failure.
Question~7.3 of~\cite{Gyoda} --- whether a generalized Markov number
determines the generalized Markov triple in which it is the largest member ---
is the direct transcription of Frobenius's uniqueness conjecture, and
\cite{Gyoda} records that it is already answered in the negative; Conjecture~7.6
is the generalization of an equivalent reformulation that Gyoda states instead,
and for which he records that no counterexample was known. What
Theorem~\ref{thm:main}(1) exhibits is a collision inside a single tree. The
nodes at $1/4$ and $2/3$ of $\GMT(22,0,5,\id)$ are distinct, namely
$((1,1),(889,3),(113,2))$ and $((15,3),(889,1),(2,2))$, yet $889$ stands in the
second of the three components at each --- the component recorded by the Farey
label --- so that $n_{1/4}=n_{2/3}=889$. Their position tags present the two
nodes as the ordered solutions $(1,113,889)$ and $(889,2,15)$
of~\eqref{eq:GM} at $(k_1,k_2,k_3)=(22,0,5)$. In each of these $889$ is the
largest entry in size, while its coordinate position is $3$ in the first and
$1$ in the second; and the ordering is not free, since of the six orderings of
$2$, $15$, $889$ only $(889,2,15)$ satisfies the equation. That is precisely
what Conjecture~7.6 forbids.

A different generalization of the uniqueness conjecture for
equation~\eqref{eq:GM} was proposed, under the name \emph{Generalized
Uniqueness Conjecture}, by Chen and Jia \cite[Conjecture~8.2]{ChenJia}, posted
on 5~November~2025 and so before~\cite{Gyoda}; Section~8.2 of~\cite{ChenJia}
develops a heuristic method for locating counterexamples to it. That
conjecture is a statement about triples ordered by size rather than about the
labeling $t\mapsto n_t$ of a fixed tree, and \eqref{eq:GM} is not symmetric in
$x,y,z$ unless $k_1=k_2=k_3$: the two solutions exhibited in Section~2 occupy
different coordinate positions, $(1,113,889)$ having its largest entry in
position~$3$ and $(889,2,15)$ in position~$1$. We make no claim about
\cite[Conjecture~8.2]{ChenJia} and do not address it here.

We have not found a previous refutation of Conjecture~7.6, nor a
previous appearance of the collision at $889$: a search of the generalized
Markov literature, including \cite{ChenJia,Gyoda,GyodaMatsushita}, turned up
the conjecture itself and the heuristic search method of
\cite[Section~8.2]{ChenJia}, but no counterexample. That search was not
exhaustive and has not been through external review, so any earlier refutation
we have missed takes precedence over this note.

\paragraph{Exact verification.}
Every assertion above is a finite computation with integers, and the data
printed here fix it completely: the four mutation evaluations of Section~2 ---
$15$, $113$ and $889$ along the left spine, and $889$ at the right child of the
$1/2$ node --- the two substitutions displayed there, and, for
Theorem~\ref{thm:main}(2), Lemma~\ref{lem:spine} together with the congruence
$F_{r+60}\equiv F_r\pmod{10}$. All arithmetic is exact integer arithmetic, each
division performed by a mutation being exact; no floating-point arithmetic and no
randomness are involved. It was recomputed once independently from that data
alone --- the six checks of Section~2, and, for each $m$ with $0\le m\le10$, the
integrality of $q_m$ and the two labels $n_{1/d_m}$ and $n_{2/3}$ obtained by
mutating $\GMT(q_m,0,0,\id)$ from its root instead of quoting
Lemma~\ref{lem:spine} --- and every check agreed with what is reported here. A
program accompanying this note carries out that recomputation, and in addition
re-derives the same facts over wider ranges of the parameters; no part of the
verification depends on it beyond what is printed above.

\begin{thebibliography}{9}

\bibitem{ChenJia}
Z.~Chen and Z.~Jia,
\emph{Tropicalization and cluster asymptotic phenomenon of generalized Markov
equations},
\href{https://arxiv.org/abs/2511.03428v2}
{arXiv:2511.03428v2 [math.NT]} (2025).

\bibitem{Gyoda}
Y.~Gyoda,
\emph{Generalized discrete Markov spectra},
\href{https://arxiv.org/abs/2512.04547v4}
{arXiv:2512.04547v4 [math.NT]} (2025), revised 2026.

\bibitem{GyodaMatsushita}
Y.~Gyoda and K.~Matsushita,
\emph{Generalization of Markov Diophantine equation via generalized cluster
algebra},
Electron. J. Combin. \textbf{30} (2023), no.~4, Paper No.~P4.10,
\href{https://doi.org/10.37236/11420}{doi:10.37236/11420}.

\end{thebibliography}

\end{document}